%% AUTO-GENERATED by hrdaya.latex_gen
%% Data version: 1.2.0  Hash: b0f93d2af26b
%% Generated: 2026-02-11T23:12:22Z
%% Regenerate with: make latex-gen

\documentclass[11pt,a4paper]{article}

\usepackage[margin=2.5cm,bottom=3cm]{geometry}
\usepackage{fontspec}
\usepackage{xeCJK}
\usepackage{xcolor}
\usepackage{longtable}
\usepackage{array}
\usepackage{enumitem}
\usepackage[perpage]{footmisc}
\usepackage{fancyhdr}
\pagestyle{fancy}
\fancyhf{}
\renewcommand{\headrulewidth}{0.4pt}
\fancyhead[L]{\small Prajñāpāramitāhṛdaya -- Tibetan Critical Edition}
\fancyhead[R]{\small\tibetanfont བཅོམ་ལྡན་འདས་མ་ཤེས་རབ་ཀྱི་ཕ་རོལ་ཏུ་ཕྱིན་པའི་སྙིང་པོ}
\fancyfoot[C]{\thepage}

\setmainfont{Noto Serif}[
  BoldFont = {Noto Serif},
  BoldFeatures = {RawFeature=+axis{wght=700}},
]
\setsansfont{Noto Sans}
\setCJKmainfont{Noto Serif CJK SC}
\setCJKsansfont{Noto Sans CJK SC}

\newfontfamily\tibetanfont{Noto Serif Tibetan}[
  BoldFont = {Noto Serif Tibetan},
  BoldFeatures = {RawFeature=+axis{wght=700}},
]
\newcommand{\tib}[1]{{\tibetanfont\XeTeXlinebreaklocale "bo" #1}}

\newcommand{\linenum}[1]{\textbf{\small #1.}}
\newcommand{\lemma}[1]{\textbf{#1}}
\newcommand{\reading}[1]{#1}
\newcommand{\wit}[1]{\textsc{#1}}
\newcommand{\gloss}[1]{{\small\color{teal}\textit{#1}}}
\newcommand{\fluent}[1]{{\small #1}}
\newcommand{\wylie}[1]{{\small\color{gray} #1}}

\begin{document}
\begin{center}
{\Huge\bfseries\tib{བཅོམ་ལྡན་འདས་མ་ཤེས་རབ་ཀྱི་ཕ་རོལ་ཏུ་ཕྱིན་པའི་སྙིང་པོ}}\\[0.5em]
{\LARGE Prajñāpāramitāhṛdaya}\\[0.3em]
{\Large Heart Sūtra}\\[1.5em]
{\large Tibetan Critical Edition}\\[0.5em]
{\normalsize Based on Degé Kangyur (Toh 21)}\\[2em]
\end{center}

\hrule
\vspace{1em}

\section*{Transmission}

The Tibetan translation was made from Sanskrit by Vimalamitra
and Rinchen Dé, revised by Gewé Lodrö and Namkha.
The Degé edition represents the long recension with frame narrative.

\vspace{1em}
\newpage

\section*{Text}

\subsection*{Section I: Opening}
\linenum{1} \tib{ཡང་དེའི་ཚེ་བྱང་ཆུབ་སེམས་དཔའ་སེམས་དཔའ་ཆེན་པོ་འཕགས་པ་སྤྱན་རས་གཟིགས་དབང་ཕྱུག་ཤེས་རབ་ཀྱི་ཕ་རོལ་ཏུ་ཕྱིན་པ་ཟབ་མོའི་སྤྱོད་པ་ཉིད་ལ་རྣམ་པར་བལྟ་ཞིང་། ཕུང་པོ་ལྔ་པོ་དེ་དག་ལ་ཡང་རང་བཞིན་གྱིས་སྟོང་པར་རྣམ་པར་བལྟའོ།}\footnote{T251 omits frame context; Tibetan preserves 'at that time' temporal marker}\\
\wylie{yang de'i tshe byang chub sems dpa' sems dpa' chen po 'phags pa spyan ras gzigs dbang phyug shes rab kyi pha rol tu phyin pa zab mo'i spyod pa nyid la rnam par blta zhing / phung po lnga po de dag la yang rang bzhin gyis stong par rnam par blta'o/}\\[0.3em]

\vspace{0.5em}
\gloss{At that time, the bodhisattva mahāsattva noble Avalokiteśvara was looking upon the profound practice of prajñāpāramitā, and he saw that the five aggregates are empty of inherent nature.}

\vspace{0.3em}
\fluent{At that time, the bodhisattva mahāsattva noble Avalokiteśvara was looking upon the profound practice of prajñāpāramitā, and he saw that the five aggregates are empty of inherent nature.}

\vspace{1em}

\subsection*{Section II: Form and Emptiness}
\linenum{2} \tib{གཟུགས་སྟོང་པའོ། སྟོང་པ་ཉིད་གཟུགས་སོ། གཟུགས་ལས་སྟོང་པ་ཉིད་གཞན་མ་ཡིན། སྟོང་པ་ཉིད་ལས་ཀྱང་གཟུགས་གཞན་མ་ཡིན་ནོ། དེ་བཞིན་དུ་ཚོར་བ་དང་། འདུ་ཤེས་དང་། འདུ་བྱེད་དང་། རྣམ་པར་ཤེས་པ་རྣམས་སྟོང་པའོ།}\\
\wylie{gzugs stong pa'o/ stong pa nyid gzugs so/ gzugs las stong pa nyid gzhan ma yin/ stong pa nyid las kyang gzugs gzhan ma yin no/ de bzhin du tshor ba dang / 'du shes dang / 'du byed dang / rnam par shes pa rnams stong pa'o/}\\[0.3em]

\vspace{0.5em}
\gloss{Form is emptiness; emptiness is form. Emptiness is not other than form; form is not other than emptiness. Likewise, feeling, perception, formation, and consciousness are empty.}

\vspace{0.3em}
\fluent{Form is emptiness; emptiness is form. Emptiness is not other than form; form is not other than emptiness. Likewise, feeling, perception, formation, and consciousness are empty.}

\vspace{1em}

\subsection*{Section III: Characteristics}
\linenum{3} \tib{ཤཱ་རིའི་བུ། དེ་ལྟར་ཆོས་ཐམས་ཅད་སྟོང་པ་ཉིད་དེ། མཚན་ཉིད་མེད་པ། མ་སྐྱེས་པ། མ་འགགས་པ། དྲི་མ་མེད་པ། དྲི་མ་དང་བྲལ་བ་མེད་པ། བྲི་བ་མེད་པ། གང་བ་མེད་པའོ།}\\
\wylie{shA ri'i bu/ de ltar chos thams cad stong pa nyid de/ mtshan nyid med pa/ ma skyes pa/ ma 'gags pa/ dri ma med pa/ dri ma dang bral ba med pa/ bri ba med pa/ gang ba med pa'o/}\\[0.3em]

\vspace{0.5em}
\gloss{Śāriputra, thus all dharmas are emptiness: without characteristics, not arising, not ceasing, not defiled, not undefiled, not decreasing, not increasing.}

\vspace{0.3em}
\fluent{Śāriputra, thus all dharmas are emptiness: without characteristics, not arising, not ceasing, not defiled, not undefiled, not decreasing, not increasing.}

\vspace{1em}

\subsection*{Section IV: Negations: Skandhas}
\linenum{4} \tib{ཤཱ་རིའི་བུ། དེ་ལྟ་བས་ན་སྟོང་པ་ཉིད་ལ་གཟུགས་མེད། ཚོར་བ་མེད། འདུ་ཤེས་མེད། འདུ་བྱེད་རྣམས་མེད། རྣམ་པར་ཤེས་པ་མེད།}\\
\wylie{shA ri'i bu/ de lta bas na stong pa nyid la gzugs med/ tshor ba med/ 'du shes med/ 'du byed rnams med/ rnam par shes pa med/}\\[0.3em]

\vspace{0.5em}
\gloss{Śāriputra, therefore in emptiness there is no form, no feeling, no perception, no formations, no consciousness;}

\vspace{0.3em}
\fluent{Śāriputra, therefore in emptiness there is no form, no feeling, no perception, no formations, no consciousness;}

\vspace{1em}

\subsection*{Section V: Negations: Āyatanas}
\linenum{5} \tib{མིག་མེད། རྣ་བ་མེད། སྣ་མེད། ལྕེ་མེད། ལུས་མེད། ཡིད་མེད། གཟུགས་མེད། སྒྲ་མེད། དྲི་མེད། རོ་མེད། རེག་བྱ་མེད། ཆོས་མེད་དོ།}\\
\wylie{mig med/ rna ba med/ sna med/ lce med/ lus med/ yid med/ gzugs med/ sgra med/ dri med/ ro med/ reg bya med/ chos med do/}\\[0.3em]

\vspace{0.5em}
\gloss{no eye, no ear, no nose, no tongue, no body, no mind; no form, no sound, no smell, no taste, no touch, no mental objects.}

\vspace{0.3em}
\fluent{no eye, no ear, no nose, no tongue, no body, no mind; no form, no sound, no smell, no taste, no touch, no mental objects.}

\vspace{1em}

\subsection*{Section VI: Negations: Dhātus}
\linenum{6} \tib{མིག་གི་ཁམས་མེད་པ་ནས་ཡིད་ཀྱི་ཁམས་མེད། ཡིད་ཀྱི་རྣམ་པར་ཤེས་པའི་ཁམས་ཀྱི་བར་དུ་ཡང་མེད་དོ།}\\
\wylie{mig gi khams med pa nas yid kyi khams med/ yid kyi rnam par shes pa'i khams kyi bar du yang med do/}\\[0.3em]

\vspace{0.5em}
\gloss{No eye-element up to no mind-element; no mind-consciousness element either.}

\vspace{0.3em}
\fluent{No eye-element up to no mind-element; no mind-consciousness element either.}

\vspace{1em}

\subsection*{Section VII: Negations: Pratītyasamutpāda}
\linenum{7} \tib{མ་རིག་པ་མེད། མ་རིག་པ་ཟད་པ་མེད་པ་ནས། རྒ་ཤི་མེད། རྒ་ཤི་ཟད་པའི་བར་དུ་ཡང་མེད་དོ།}\footnote{Tibetan uses zad pa (exhaustion) matching Sanskrit kṣaya rather than standard nirodha - evidence for Nattier's back-translation thesis}\\
\wylie{ma rig pa med/ ma rig pa zad pa med pa nas/ rga shi med/ rga shi zad pa'i bar du yang med do/}\\[0.3em]

\vspace{0.5em}
\gloss{No ignorance, no ending of ignorance, up to no old age and death, no ending of old age and death.}

\vspace{0.3em}
\fluent{No ignorance, no ending of ignorance, up to no old age and death, no ending of old age and death.}

\vspace{1em}

\subsection*{Section VIII: Negations: Noble Truths}
\linenum{8} \tib{སྡུག་བསྔལ་བ་དང་། ཀུན་འབྱུང་བ་དང་། འགོག་པ་དང་། ལམ་མེད། ཡེ་ཤེས་མེད། ཐོབ་པ་མེད། མ་ཐོབ་པ་ཡང་མེད་དོ།}\\
\wylie{sdug bsngal ba dang / kun 'byung ba dang / 'gog pa dang / lam med/ ye shes med/ thob pa med/ ma thob pa yang med do/}\\[0.3em]

\vspace{0.5em}
\gloss{No suffering, no origination, no cessation, no path; no wisdom, no attainment, and no non-attainment.}

\vspace{0.3em}
\fluent{No suffering, no origination, no cessation, no path; no wisdom, no attainment, and no non-attainment.}

\vspace{1em}

\subsection*{Section IX: Bodhisattva Result}
\linenum{9} \tib{ཤཱ་རིའི་བུ། དེ་ལྟ་བས་ན་བྱང་ཆུབ་སེམས་དཔའ་རྣམས་ཐོབ་པ་མེད་པའི་ཕྱིར། ཤེས་རབ་ཀྱི་ཕ་རོལ་ཏུ་ཕྱིན་པ་ལ་བརྟེན་ཅིང་གནས་ཏེ། སེམས་ལ་སྒྲིབ་པ་མེད་པས་འཇིགས་པ་མེད་དེ། ཕྱིན་ཅི་ལོག་ལས་ཤིན་ཏུ་འདས་ནས་མྱ་ངན་ལས་འདས་པའི་མཐར་ཕྱིན་ཏོ།}\\
\wylie{shA ri'i bu/ de lta bas na byang chub sems dpa' rnams thob pa med pa'i phyir/ shes rab kyi pha rol tu phyin pa la brten cing gnas te/ sems la sgrib pa med pas 'jigs pa med de/ phyin ci log las shin tu 'das nas mya ngan las 'das pa'i mthar phyin to/}\\[0.3em]

\vspace{0.5em}
\gloss{Śāriputra, therefore, because there is no attainment, bodhisattvas rely on and abide in prajñāpāramitā. Because their minds have no obscurations, they have no fear. Having completely transcended error, they attain the final nirvāṇa.}

\vspace{0.3em}
\fluent{Śāriputra, therefore, because there is no attainment, bodhisattvas rely on and abide in prajñāpāramitā. Because their minds have no obscurations, they have no fear. Having completely transcended error, they attain the final nirvāṇa.}

\vspace{1em}

\subsection*{Section X: Buddha Result}
\linenum{10} \tib{དུས་གསུམ་དུ་རྣམ་པར་བཞུགས་པའི་སངས་རྒྱས་ཐམས་ཅད་ཀྱང་ཤེས་རབ་ཀྱི་ཕ་རོལ་ཏུ་ཕྱིན་པ་ལ་བརྟེན་ནས་བླ་ན་མེད་པ་ཡང་དག་པར་རྫོགས་པའི་བྱང་ཆུབ་ཏུ་མངོན་པར་རྫོགས་པར་སངས་རྒྱས་སོ།}\\
\wylie{dus gsum du rnam par bzhugs pa'i sangs rgyas thams cad kyang shes rab kyi pha rol tu phyin pa la brten nas bla na med pa yang dag par rdzogs pa'i byang chub tu mngon par rdzogs par sangs rgyas so/}\\[0.3em]

\vspace{0.5em}
\gloss{All the buddhas who abide in the three times, relying on prajñāpāramitā, fully awaken to unsurpassed, perfect, complete enlightenment.}

\vspace{0.3em}
\fluent{All the buddhas who abide in the three times, relying on prajñāpāramitā, fully awaken to unsurpassed, perfect, complete enlightenment.}

\vspace{1em}

\subsection*{Section XI: Mantra Praise}
\linenum{11} \tib{དེ་ལྟ་བས་ན་ཤེས་རབ་ཀྱི་ཕ་རོལ་ཏུ་ཕྱིན་པའི་སྔགས། རིག་པ་ཆེན་པོའི་སྔགས། བླ་ན་མེད་པའི་སྔགས། མི་མཉམ་པ་དང་མཉམ་པའི་སྔགས། སྡུག་བསྔལ་ཐམས་ཅད་རབ་ཏུ་ཞི་བར་བྱེད་པའི་སྔགས། མི་རྫུན་པས་ན་བདེན་པར་ཤེས་པར་བྱ་སྟེ། ཤེས་རབ་ཀྱི་ཕ་རོལ་ཏུ་ཕྱིན་པའི་སྔགས་སྨྲས་པ།}\\
\wylie{de lta bas na shes rab kyi pha rol tu phyin pa'i sngags/ rig pa chen po'i sngags/ bla na med pa'i sngags/ mi mnyam pa dang mnyam pa'i sngags/ sdug bsngal thams cad rab tu zhi bar byed pa'i sngags/ mi rdzun pas na bden par shes par bya ste/ shes rab kyi pha rol tu phyin pa'i sngags smras pa/}\\[0.3em]

\vspace{0.5em}
\gloss{Therefore, the prajñāpāramitā mantra, the mantra of great knowledge, the unsurpassed mantra, the mantra equal to the unequalled, the mantra that completely pacifies all suffering, should be known as true because it is not false. The prajñāpāramitā mantra is proclaimed:}

\vspace{0.3em}
\fluent{Therefore, the prajñāpāramitā mantra, the mantra of great knowledge, the unsurpassed mantra, the mantra equal to the unequalled, the mantra that completely pacifies all suffering, should be known as true because it is not false. The prajñāpāramitā mantra is proclaimed:}

\vspace{1em}

\subsection*{Section XII: Mantra}
\linenum{12} \tib{ཏདྱ་ཐཱ། ཨོཾ་ག་ཏེ་ག་ཏེ་པཱ་ར་ག་ཏེ་པཱ་ར་སཾ་ག་ཏེ་བོ་དྷི་སྭཱ་ཧཱ།}\\
\wylie{tad+ya thA/ oM ga te ga te pA ra ga te pA ra saM ga te bo d+hi swA hA/}\\[0.3em]

\vspace{0.5em}
\gloss{Tadyathā: Oṃ gate gate pāragate pārasaṃgate bodhi svāhā.}

\vspace{0.3em}
\fluent{Tadyathā: Oṃ gate gate pāragate pārasaṃgate bodhi svāhā.}

\vspace{1em}

\end{document}
